\section{Introduction}

Forecasting under partial observation is not only a problem of minimizing a
single error column. In influenza hospitalization forecasting, the observed
target is a reported hospitalization-rate time series, while many candidate
models reason in terms of latent quantities such as infectious prevalence,
hospitalized compartments, or delayed observation processes. A model may appear
weak because the latent dynamics are wrong, because the observation map is
wrong, because a simple non-mechanistic baseline is sufficient, or because the
evidence being used does not match the claim being made.

This creates a recommendation problem rather than a pure leaderboard problem.
The central question is:
\begin{quote}
How can agreement and disagreement across model families, proposers, policies,
and evidence modes be used to decide which forecasting recommendations are
credible under partial observation, while preventing leakage and overclaiming?
\end{quote}

We answer this question with a verifier-gated model-selection framework
(Figure~\ref{fig:framework}). The framework treats every model suggestion as a
structured candidate record. Candidates may come from deterministic heuristics,
random orderings, failure-guided refinement, or constrained LLM-backed
proposers, but all candidates enter the same verifier. The verifier enforces
schema validity, explicit allowlists, no-code-generation, no-leakage, budget,
observation-label, and claim-safety constraints. Only verified candidates are
evaluated by deterministic code. Selection policies operate on validation or
rolling-validation evidence, while held-out test metrics are attached only
after selection as post-hoc descriptive evidence.

\begin{figure*}[t]
\centering
\includegraphics[width=\textwidth]{fig1_agentic_framework_pipeline_fancy.pdf}
\caption{Verifier-gated agentic framework for credible cross-family
model selection under partial observation. The constrained agentic proposer
emits structured candidate records from an explicit cross-family hypothesis
space. A central verifier gate enforces schema validity, allowlist membership,
observation compatibility, budget constraints, no-leakage evidence use, no
code generation, and claim safety before any candidate reaches the evidence
engines. Evidence-aware policies and a claim auditor convert non-test evidence
into an audited recommendation rather than a universal-winner claim.}
\label{fig:framework}
\end{figure*}

The contribution is methodological, not a claim of operational flu forecasting
superiority. We make six contributions:
\begin{enumerate}
\item a verifier-gated cross-family model-selection framework for partially
observed forecasting;
\item a constrained agentic proposer protocol with task schemas, explicit
allowlists, JSON candidate records, and verifier feedback;
\item a no-leakage evidence protocol with Pareto-epsilon, weighted-score, and
hard-veto policies plus claim-boundary auditing;
\item synthetic structured recovery experiments that test observation-label
and delay-label recovery under known ground truth;
\item candidate-budget replay and bounded real-data execution experiments that
evaluate proposer ordering without unrestricted refitting; and
\item a FluSurv-NET case study and four-provider agentic proposer benchmark over
OpenAI/GPT, Claude, Gemini, and DeepSeek under the same verifier protocol.
\end{enumerate}

\section{Related Work}

Operational influenza forecasting has established the value of collaborative
multi-model assessment, real-time evaluation, and ensemble forecasts
\citep{shaman2012forecasting,shaman2013realtime,reich2019collaborative,
reich2019accuracy,lutz2019pathforward,ray2024flusion}. Our setting is
different. We use FluSurv-NET as a retrospective case study for credible
model recommendation under partial observation, not as a FluSight leaderboard
entry or operational forecasting system.

Mechanistic epidemic forecasting and symbolic discovery provide interpretable
structure, but they also introduce model-space and observation-map choices.
SINDy, AI Feynman, and PySR search for interpretable equations from data
\citep{brunton2016sindy,udrescu2020aifeynman,cranmer2023pysr}. Our structured
search is deliberately narrower: the compartment templates, observation maps,
and delays are bounded in advance, and all discovered candidates are compared
against simple forecasting baselines, hand-specified mechanistic baselines, and
same-grammar ablations.

Recent language-model systems for scientific search motivate proposal and
evaluation loops \citep{romera2024funsearch,shojaee2024llmsr,lu2024aiscientist}.
We adopt the loop structure but restrict the autonomy. A language model, when
used, is a structured proposer. It cannot invent model families, generate
executable model code, change the allowlist, execute forecasts, or view
held-out test metrics.

Adjacent AI-for-science pipelines reinforce the same separation between
proposal, verification, evaluation, and claim discipline. In molecular design,
REINVENT-style workflows, scaffold-constrained DrugEx search, multi-target
molecular generation, and pocket-conditioned diffusion propose or optimize
candidate objects, but their scientific value depends on downstream scoring,
filtering, multi-objective evaluation, uncertainty, and domain-specific
validation rather than on generation alone
\citep{loeffler2024reinvent,liu2023drugex,ai2024mtmolgpt,lin2024diffbp}.
Platform and benchmarking work makes this boundary explicit: Chemistry42-style
optimization and MolScore-style evaluation treat proposal quality as something
that must be filtered, scored, and benchmarked under explicit objectives, while
uncertainty and conformal-prediction studies emphasize applicability-domain and
confidence-aware reporting \citep{ivanenkov2023chemistry42,thomas2024molscore,
dutschmann2023uncertainty,arvidsson2024cpsign}. Prospective de novo design and
mechanism-aware foundation-model discussions further motivate closed-loop
feedback and mechanism-grounded evaluation, but they also illustrate why a
proposal engine should not be confused with a validated scientific conclusion
\citep{atz2024prospective,lavecchia2026mechanismaware}. We use these
cross-domain references only to motivate proposer/evaluator separation,
feedback, uncertainty, and claim boundaries; they do not support influenza
forecasting performance, FluSurv-NET recommendations, provider superiority, or
real-world epidemic mechanism recovery.

This design follows standard concerns in out-of-sample and rolling-origin
forecast evaluation \citep{tashman2000outofsample,hyndman2021fpp3} and uses
multi-objective selection ideas related to Pareto optimization
\citep{deb2002nsga2}.
